%!TeX program = xelatex
\documentclass{SYSUReport}
\usepackage{tabularx}
\usepackage{float}
% 根据个人情况修改
\headl{}
\headc{}
\headr{并行程序设计与算法实验}
\lessonTitle{并行程序设计与算法实验}
\reportTitle{Lab6-Pthreads并行构造}
\stuname{林宏宇}
\stuid{23320093}
\inst{计算机学院}
\major{计算机科学与技术}
\date{2025年5月18日}

\begin{document}

% =============================================
% Part 1: 封面
% =============================================
\cover
\thispagestyle{empty}
\clearpage

\section{实验目的}
\begin{itemize}
  \item 深入理解C/C++程序编译、链接与构建过程
  \item 提高Pthreads多线程编程综合能力
  \item 并行库开发与性能验证
\end{itemize}

\section{实验内容}
\begin{itemize}
  \item 基于Pthreads的并行计算库实现：parallel\_for + 调度策略
  \item 动态链接库生成和使用
  \item 典型问题的并行化改造：矩阵乘法、热传导问题
\end{itemize}

\section{实验过程}

\subsection{环境与工具}
简要说明实验所使用的操作系统、编译器 (gcc/g++) 版本、以及 Pthreads 库。
\begin{itemize}
    \item 操作系统: Ubuntu 22.04 (Linux 6.11.0-25-generic, x86\_64)
    \item 编译器: GCC 11.4.0
    \item CPU: Intel Xeon Gold 6240R @ 2.40GHz, 24核/48线程 (双路共96逻辑核)
    \item Pthreads: Linux glibc 自带线程库 (NPTL), 编译时添加 \texttt{-pthread} 选项
\end{itemize}

\subsection{核心函数实现 (\texttt{parallel\_for})}
简要描述 \texttt{parallel\_for} 函数的关键设计思路和实现。重点说明线程创建、任务划分和同步机制。

\textbf{设计思路：}
\texttt{parallel\_for} 函数模仿 OpenMP 的并行 for 循环机制，基于 Pthreads 实现。核心设计包括：

\begin{enumerate}
    \item \textbf{持久线程池：} 避免每次调用都创建/销毁线程的开销。首次调用时创建指定数量的工作线程，线程通过条件变量等待任务分配。使用代际计数器 (generation counter) 区分不同批次的并行任务，避免工作线程在 cleanup 阶段与新任务之间的竞态条件。
    \item \textbf{静态调度 (Static)：} 将迭代空间 $[start, end)$ 按线程数均匀分块。前 $remainder$ 个线程多分配一个迭代，其余线程各分配 $total/num\_threads$ 个迭代。每个线程处理连续的迭代范围，适合每次迭代工作量相近的场景。
    \item \textbf{动态调度 (Dynamic)：} 使用互斥锁保护的共享计数器 \texttt{w\_next\_iter} 实现任务窃取。每个线程循环获取 \texttt{chunk\_size} 个迭代，直到所有迭代被分配完毕。适合每次迭代工作量不均衡的场景，但额外的锁竞争带来开销。
    \item \textbf{同步机制：} 主线程通过 \texttt{pthread\_cond\_broadcast} 唤醒所有工作线程，通过 \texttt{pthread\_cond\_wait} 等待所有工作线程完成。工作线程完成工作后递增完成计数器，最后一个完成的线程通知主线程。
\end{enumerate}

\textbf{函数接口：}
\begin{verbatim}
void parallel_for(int start, int end, int inc,
                  void *(*functor)(int, void*), void *arg,
                  int num_threads);

void parallel_for_sched(int start, int end, int inc,
                        void *(*functor)(int, void*), void *arg,
                        int num_threads, int schedule, int chunk_size);
\end{verbatim}

其中 \texttt{schedule} 取 \texttt{SCHED\_STATIC (0)} 或 \texttt{SCHED\_DYNAMIC (1)}，\texttt{chunk\_size} 指定动态调度的块大小。

\subsection{动态库生成与使用}
说明生成动态链接库 (.so) 的主要命令或 Makefile 规则，并简述如何在主程序 (如矩阵乘法、热传导) 中链接和调用该库。

\textbf{生成动态库：}
\begin{verbatim}
# 编译位置无关代码
gcc -Wall -O2 -fPIC -c parallel_for.c
# 创建动态库
gcc -shared -Wl,-soname,libparallel_for.so.1 \
    -o libparallel_for.so.1.0 parallel_for.o -pthread
# 创建符号链接
ln -sf libparallel_for.so.1.0 libparallel_for.so.1
ln -sf libparallel_for.so.1 libparallel_for.so
\end{verbatim}

\textbf{链接主程序：}
\begin{verbatim}
# 编译主程序并链接自定义动态库
gcc -Wall -O2 -c heated_plate_pthreads.c
gcc -o heated_plate_pthreads heated_plate_pthreads.o \
    -L. -lparallel_for -Wl,-rpath,. -pthread -lm
\end{verbatim}

其中 \texttt{-L.} 指定当前目录为库搜索路径，\texttt{-lparallel\_for} 链接 \texttt{libparallel\_for.so}，\texttt{-Wl,-rpath,.} 设置运行时库搜索路径。

\subsection{应用测试 (矩阵乘法)}
为验证 \texttt{parallel\_for} 的正确性与效率，首先使用矩阵乘法进行测试。测试矩阵大小为 512×512，使用 4 线程静态调度。

\textbf{正确性验证：} 将并行计算结果与串行参考结果进行逐元素比较，最大误差在机器精度范围内 ($<10^{-6}$)，验证了 \texttt{parallel\_for} 库任务划分和同步机制的正确性。矩阵乘法同时检验了 functor 接口对二维循环的适配能力——将内层循环封装为按行处理的 functor，每行完成一整列的乘加操作。

\textbf{效率验证：} 矩阵乘法的时间复杂度为 $O(N^3)$，计算密度高（每个输出元素需 $2N$ 次浮点运算），是验证并行效率的理想负载。4 线程静态调度下，矩阵规模 $N=512$ 时总计算量约 2.7 亿次浮点运算，并行版本在 4 线程下获得了接近线性加速的效果。该测试表明 \texttt{parallel\_for} 在计算密集型任务上能够有效利用多核资源，线程创建和同步的开销相对于计算量可忽略不计。

\subsection{应用测试 (热传导)}
简述如何将 \texttt{parallel\_for} 应用于热传导问题，替换原有的并行机制。

\begin{itemize}
    \item \textbf{改造方法：} 原 OpenMP 实现使用 \texttt{\#pragma omp parallel for} 指令对热传导迭代的各阶段进行并行化。Pthreads 改造版本将每个并行循环封装为独立的 functor 函数，通过 \texttt{parallel\_for} 调用。定义统一的上下文结构体 \texttt{hp\_ctx} 包含网格数组指针 (w, u)、网格尺寸、均值、差值及互斥锁。
    \item \textbf{Functor 设计：} 对于边界条件设置等一维循环，functor 接收行索引 i 或列索引 j 直接操作；对于内点计算等二维循环，functor 接收行索引 i，内部对整行的 j 循环进行处理，以减少 parallel\_for 的函数指针调用开销。均值求和的归约操作通过互斥锁保护实现。
    \item \textbf{关键优化——diff\_row 去锁：} 最大差值归约的初始实现对每一行都执行 \texttt{pthread\_mutex\_lock/unlock} 更新全局最大值，每轮迭代产生约 500 次锁操作，16955 轮累计超 800 万次。实测发现该实现下 4 线程及以上几乎不加速（4 线程 8.08s vs 2 线程 8.64s），因为锁竞争随线程数增加反而恶化。优化方案为：在上下文结构体中增加 \texttt{row\_diff[M]} 数组，每个 functor 将本行局部最大值直接写入 \texttt{row\_diff[i]}（无锁操作，各行独立），待 \texttt{parallel\_for} 返回后由主线程串行遍历 $[1, M-1]$ 取全局最大值——串行遍历 500 个元素的代价可忽略不计。该优化将 4 线程耗时从 8.08s 降至 4.72s（提升 41\%），是本次实验中性能提升最大的单项改进。
    \item \textbf{测试设置：} 网格大小 500×500，收敛精度 $\epsilon = 0.001$。线程数取 1、2、4、8、16。调度方式分别测试静态调度 (Static) 和动态调度 (Dynamic, chunk=1, 4, 16)。
\end{itemize}

\newpage
\section{实验结果与分析}

\subsection{性能测试结果}
展示不同线程数和调度方式下，自定义 Pthreads 实现与原始 OpenMP 实现的性能对比。

\begin{table}[h!]
\centering
\caption{热传导问题性能对比 (Pthreads vs OpenMP, 网格大小: 500×500, 收敛精度 $\epsilon=0.001$)}
\begin{tabular}{|c|c|c|c|c|}
\hline
\textbf{线程数} & \textbf{调度方式} & \textbf{Pthreads} & \textbf{OpenMP} & \textbf{Pthreads/OpenMP} \\
 &  & 时间 (s) & 时间 (s) & 时间比 \\
\hline
1 (串行) & N/A & 14.51 & 13.63 & 1.06 \\
\hline
\multicolumn{5}{|c|}{Pthreads: 静态调度 (Static)} \\
\hline
2 & Static & 8.70 & 6.94 & 1.25 \\
4 & Static & 4.72 & 3.63 & 1.30 \\
8 & Static & 4.54 & 2.11 & 2.15 \\
16 & Static & 4.92 & 1.44 & 3.42 \\
\hline
\multicolumn{5}{|c|}{Pthreads: 动态调度 (Dynamic, chunk=16)} \\
\hline
2 & Dynamic,16 & 8.12 & 6.94 & 1.17 \\
4 & Dynamic,16 & 4.97 & 3.63 & 1.37 \\
8 & Dynamic,16 & 5.14 & 2.11 & 2.43 \\
16 & Dynamic,16 & 3.93 & 1.44 & 2.73 \\
\hline
\multicolumn{5}{|c|}{Pthreads: 动态调度 (Dynamic, chunk=4)} \\
\hline
2 & Dynamic,4 & 8.88 & 6.94 & 1.28 \\
4 & Dynamic,4 & 5.64 & 3.63 & 1.55 \\
8 & Dynamic,4 & 4.88 & 2.11 & 2.31 \\
16 & Dynamic,4 & 7.37 & 1.44 & 5.12 \\
\hline
\multicolumn{5}{|c|}{Pthreads: 动态调度 (Dynamic, chunk=1)} \\
\hline
2 & Dynamic,1 & 12.65 & 6.94 & 1.82 \\
4 & Dynamic,1 & 10.11 & 3.63 & 2.79 \\
8 & Dynamic,1 & 12.87 & 2.11 & 6.09 \\
16 & Dynamic,1 & 13.08 & 1.44 & 9.09 \\
\hline
\end{tabular}
\label{tab:performance_comparison}
\end{table}

\begin{table}[h!]
\centering
\caption{加速比对比 (相对于各自单线程基线)}
\begin{tabular}{|c|c|c|c|c|c|}
\hline
\textbf{线程数} & \textbf{OpenMP} & \textbf{Static} & \textbf{Dyn,16} & \textbf{Dyn,4} & \textbf{Dyn,1} \\
\hline
1 & 1.00× & 1.00× & 1.00× & 1.00× & 1.00× \\
2 & 1.96× & 1.67× & 1.82× & 1.66× & 1.19× \\
4 & 3.76× & 3.08× & 2.98× & 2.62× & 1.48× \\
8 & 6.45× & 3.19× & 2.88× & 3.03× & 1.17× \\
16 & 9.49× & 2.95× & 3.77× & 2.01× & 1.15× \\
\hline
\end{tabular}
\label{tab:speedup}
\end{table}

\begin{table}[h!]
\centering
\caption{动态调度不同 chunk 大小在 4 线程下的性能对比}
\begin{tabular}{|c|c|c|c|}
\hline
\textbf{chunk大小} & \textbf{时间 (s)} & \textbf{加速比} & \textbf{相对Static} \\
\hline
Static (基线) & 4.72 & 3.08× & -- \\
1 & 10.11 & 1.48× & 慢 114\% \\
4 & 5.64 & 2.62× & 慢 19\% \\
16 & 4.97 & 2.98× & 慢 5\% \\
\hline
\end{tabular}
\label{tab:chunk_comparison}
\end{table}

所有配置（OpenMP 与 Pthreads 各调度策略/线程数）均在同一网格规模 (500×500) 和收敛精度 ($\epsilon=0.001$) 下收敛于 16955 次迭代，均值均为 74.949900。迭代次数和最终结果完全一致，验证了 Pthreads 改造版本计算逻辑的正确性。

\newpage
\subsection{结果分析与总结}

\textbf{1. 正确性：} 所有配置在相同条件下均收敛于 16955 次迭代，最终均值 74.949900 完全一致，证明 parallel\_for 库的改造未引入计算错误。

\textbf{2. 并行加速效果：} Pthreads 实现成功实现了并行加速，但效果因线程数和调度策略差异显著：
\begin{itemize}
    \item \textbf{静态调度}在 4 线程加速良好 (3.08×)，但 8 线程后出现严重扩展瓶颈（8 线程 3.19×，边际收益接近于零），16 线程因同步开销和 NUMA 效应反而退化为 2.95×。原因是线程池同步开销随线程数增加而增大：每次 parallel\_for 调用需广播唤醒所有池中线程并等待全部回报，即使部分线程不参与计算也需同步。此外，双路 NUMA 架构下跨 socket 访存延迟进一步限制了高线程数的扩展性。
    \item \textbf{动态调度 chunk=16} 在 16 线程时达到 pthreads 最优加速比 (3.77×)，且超过同线程数静态调度 (2.95×)。因为较大 chunk 降低了锁获取频率，同时动态调度使得与数据在同一 socket 的"近端"线程能抢占更多工作，部分弥补了 NUMA 效应。
    \item \textbf{动态调度 chunk=1} 几乎无法加速，8 线程仅为 1.17×。每获取一行就需要一次加锁/解锁，498 行 × 3 次 parallel\_for × 16955 迭代产生约 2500 万次锁操作，竞争完全抵消并行收益。
\end{itemize}

\textbf{3. 与 OpenMP 对比：} Pthreads 在所有配置下均慢于 OpenMP，且差距随线程数扩大 (2线程 1.25× → 16线程 3.42×)。主要原因：
\begin{itemize}
    \item \textbf{编译器优化差异：} OpenMP 由编译器将循环体直接编译为内联代码，而 Pthreads 通过函数指针调用 functor，每次调用有间接跳转开销。
    \item \textbf{运行时开销：} OpenMP 的并行区域保持线程存活，内部多个 \texttt{\#pragma omp for} 之间仅需 barrier 同步。Pthreads 每次 parallel\_for 调用都是一次完整的 fork/join 周期（broadcast→work→gather），热传导每轮迭代 3 次调用，16955 轮共计约 5 万次同步周期。
    \item \textbf{NUMA 感知：} OpenMP 运行时可能采用 first-touch 策略在并行循环初始化时分散内存分配，而 Pthreads 版本由主线程统一 malloc，所有数据集中在一个 socket 的内存上，另一 socket 线程远程访存延迟更高。
\end{itemize}

\textbf{4. 调度策略比较：}
\begin{itemize}
    \item 对于热传导这种每次迭代工作量均匀的应用，静态调度应是理论最优。但实验中动态调度 chunk=16 在 16 线程时反超静态 (3.93s vs 4.92s)，说明静态调度对 NUMA 效应更敏感。
    \item 增大 chunk 大小显著改善动态调度性能：4 线程下 chunk=16 仅比静态慢 5\%，而 chunk=1 慢 114\%。
\end{itemize}

\end{document}
